\section{Related Work}
\paragraph{AI for Research.}
Many works have studied how to use LMs for different components of the scientific research pipeline, including literature search \citep{openai_deep_research, zheng-etal-2025-deepresearcher}, idea generation \citep{si2024can, wang-etal-2024-scimon}, and idea execution \citep{si2025ideation, si2026towards, karpathy2026autoresearch}. A concurrent work trains models to judge research impact from citation signals, then uses the learned judge as a reward model for idea generation \citep{tong2026ai}. In contrast, we focus on the problem of synthesizing insights from two parent papers, using similarity to ground-truth downstream insights as a more direct training signal. A recent line of work studies scientific progress as a forecasting problem. Most closely related concurrent work introduces \textsc{PreScience}, a benchmark that decomposes the research process into collaborator prediction, prior work selection, contribution generation, and impact prediction \citep{ajith2026prescience}. In their contribution generation task, models generate a future paper's title and abstract conditioned on prior work and other historical context, and are evaluated using an LM-based similarity metric. Our work is complementary but more focused: rather than modeling the full scientific workflow, we isolate the idea-synthesis problem and ask whether a model can generate the core insight of a downstream paper from a given pair of parent papers. Thus, while \textsc{PreScience} studies broad scientific forecasting, we study a more controlled and direct test of whether models can synthesize a novel downstream insight from specific prior works.


\paragraph{Literature-based Discovery.}
Literature-Based Discovery (LBD) is a field that involves using computational methods to uncover hidden links between seemingly unrelated bodies of research, aiming to infer novel and interesting knowledge \citep{swanson1986fish, swanson2008literature}. For instance, \cite{swanson1986undiscovered} noticed that papers on fish oil mentioned it reduces blood viscosity, while others linked high blood viscosity to Raynaud’s syndrome. Therefore, \cite{swanson1986undiscovered} hypothesized that fish oil could treat Raynaud’s, a prediction later confirmed by clinical research. However, existing LBD approaches struggle to scale to the massive volume of today’s literature and often rely on domain-specific knowledge sources. Moreover, they typically output only ranked lists of candidate connections rather than clear, hypothesis-like insights that researchers can act on \citep{sebastian2017emerging, ganiz2005recent, bekhuis2006conceptual, smalheiser2012literature}. 

% Scientific discovery often depends on synthesizing information across diverse sources to generate novel hypotheses or insights that are not explicitly stated in any single document \citep{swanson1991complementary}. 